\documentclass[11pt,a4paper]{article}
\usepackage[utf8]{inputenc}
\usepackage[T1]{fontenc}
\usepackage{amsmath,amssymb,amsthm}
\usepackage{geometry}
\usepackage{xcolor}
\usepackage{hyperref}
\geometry{margin=2.5cm}
\hypersetup{colorlinks=true,linkcolor=blue!60!black,citecolor=blue!60!black,urlcolor=blue!60!black}
\newtheorem{theorem}{Theorem}
\newtheorem{remark}[theorem]{Remark}
\theoremstyle{definition}
\newtheorem{observation}[theorem]{Observation}

\title{\textbf{Weighing the obstruction: the smoothness of $p-1$, its parity
signature, and a measured Tamagawa reconstruction of the twin-prime constant}}
\author{Fouad Bensmail\\ Independent researcher, Constantine, Algeria\\
ORCID: 0009-0006-9541-4207}
\date{14 September 2026}

\begin{document}
\maketitle

\begin{abstract}
This paper is experimental in the sense of the journal's manifesto: every
claim is a high-precision measurement, reproducible by a single
dependency-free Python command, and each measurement carries its own guard
(a predicted ratio, a sigma budget, and a contraction test across two
scales). Four results. (1) The naive Dickman model for $P^+(p-1)$ is rejected
at $10^6$ (ratios $1.1998/1.7412/4.0423$ against $\rho(u)$); a corrected null
model (parity plus Dirichlet densities $1/(q-1)$) is then forged and weighed
at $10^6$ and $10^7$: all ratios lie within $0.2\,\sigma$ and the residual
contracts as $1/\sqrt N$, excluding a $1/\ln X$ bias. (2) Conditioning on
$p \bmod 4$, the local obstruction at $q=2$ proves computable and is captured
class by class (8/8 ratios within $2\sigma$ at both scales), while the
unconditioned model misses each class in the predicted direction
($0.63$ versus $2.42$ at $u=5$, $10^7$): structural bias, not fluctuation.
(3) The singular series of pair motifs is a relative invariant: exact under
translations and reflections, covariant under dilation $d$ by the
multiplicative character $\chi_H(d)=\prod_{p\mid d}(1-1/p)/(1-\nu_p/p)$,
verified to twelve digits; the covariance descends into prime-pair counts at
$10^7$ and $10^8$ (reproducing the published twin counts $58\,980$ and
$440\,312$ exactly). (4) Chinese Remainder independence is measured rather
than assumed: for eight moduli $D\le46\,189$ the brute-force period count
equals the product of local counts exactly, and the global frequency over
$n\le10^7$ matches the product of local frequencies to $7.11\times10^{-7}$;
the local Tamagawa volumes $\beta_p=1-1/(p-1)^2$ (odd $p$) reconstruct
$2C_2=1.3203236$ to $1.2\times10^{-7}$. The paper proves no open conjecture
and claims none: two initial predictions were corrected by the measurement
itself, and one self-detected coding error is published as a refused run.
\end{abstract}

\section{Method: every measurement carries its own guard}
A measurement without a guard is an anecdote. Each weighing below announces
beforehand what it expects (a ratio, a sigma budget), and each is run at two
scales: a deviation that contracts as $1/\sqrt N$ is a fluctuation; one that
persists is a structure. All companions are pure Python, archived with the
lab notebook on Zenodo (concept DOI 10.5281/zenodo.22314057; cahier v9:
10.5281/zenodo.22757396).

\section{The naive model and its refusal}
\begin{observation}[Dickman rejected]
At $X=10^6$, the measured proportions of $p-1$ with $P^+(p-1)\le(p-1)^{1/u}$
for $u=2,3,4$ are $0.36817/0.08464/0.02009$ against
$\rho(u)=0.30685/0.04861/0.00497$: ratios $1.1998/1.7412/4.0423$, bias
growing with $u$. The model ``$p-1$ is a random integer'' is refused, and the
refusal is published.
\end{observation}

\section{The corrected null model at two scales}
The corrected model draws $m$ even with $m+1$ free of prime factors $\le Q$,
reproducing exactly the two biases: parity, and Dirichlet densities
$P(q\mid p-1)=1/(q-1)$.
\begin{observation}[Residual evaporates]
At $10^7$ the measured proportions for $u=2,3,4,5$ are
$0.35583/0.07519/0.01449/0.00360$ against model
$0.35577/0.07513/0.01445/0.00359$: ratios $0.9998/0.9993/0.9976/0.9965$ with
sigma budgets $0.07/0.11/0.17/0.12$; at $10^6$ the ratios are
$0.9994/0.9973/0.9939$ for $u=2,3,4$. From $10^6$ to $10^7$ the deviations
contract by a factor $\approx 3$, the $1/\sqrt N$ of a fluctuation --- not
the $\times 0.86$ of a $1/\ln X$ bias.
\end{observation}

\section{The parity signature of the local obstruction at $q=2$}
\begin{observation}[Class by class]
Weighing $p\equiv1\pmod4$ ($n_1=39\,175$ then $332\,180$) and
$p\equiv3\pmod4$ ($n_3=39\,321$ then $332\,397$) separately, the gaps
$\rho_1-\rho_3$ at $u=2,3,4,5$ are
$+0.11582/+0.05657/+0.02002/+0.00738$ at $10^6$ and
$+0.09862/+0.04496/+0.01324/+0.00423$ at $10^7$; the conditioned ratios
mod1/mes1 are $0.9998/0.9992/0.9974/0.9964$ and mod3/mes3 are
$0.9999/0.9994/0.9980/0.9966$; the blind (unconditioned) ratios are
$0.8781/0.7692/0.6847/0.6276$ in class 1 and
$1.1606/1.4252/1.8360/2.4154$ in class 3. The signature $\rho_1>\rho_3$
holds 4/4 at both scales; the eight conditioned ratios lie within $2\sigma$
(8/8 and 8/8); the blind model misses each class in the predicted direction
and does not evaporate. Guards: mean $v_2(p-1)=3.00/1.00$ exact; Chebyshev
bias $+146$ then $+217$.
\end{observation}

\section{Isometries and homotheties of pair motifs}
\begin{theorem}[Relative invariance, exact]
For a motif $H$ with $\nu_p(H)$ forbidden classes mod $p$, the singular
series $\mathfrak S(H)=\prod_p (1-\nu_p/p)(1-1/p)^{-k}$ satisfies:
$\mathfrak S(H+t)=\mathfrak S(H)$ and $\mathfrak S(-H)=\mathfrak S(H)$ exactly
(per prime); and $\mathfrak S(dH)=\mathfrak S(H)\prod_{p\mid d}
(1-1/p)/(1-\nu_p/p)$. Measured/predicted $=1.000000000000$ for
$d=2,3,4,6,12$; translation and reflection deviations $=0$.
\end{theorem}
\begin{observation}[Covariance descends into counts]
At $10^7$: $C_2=58\,980$, $C_4=58\,622$, $C_6=117\,207$;
$C_4/C_2=0.9939$ ($1.04\sigma$), $C_6/C_2=1.9872$ ($2.53\sigma$).
At $10^8$: $C_2=440\,312$, $C_4=440\,258$, $C_6=879\,908$;
$C_4/C_2=0.9999$ ($0.06\sigma$), $C_6/C_2=1.9984$ ($0.88\sigma$): the
$10^7$ deviation contracts, and $C_2$ matches published twin counts exactly
at both scales.
\end{observation}

\section{Measured Tamagawa volumes and CRT independence}
\begin{observation}[Global equals product of locals]
For eight modulus sets with $D=\prod p\le46\,189$, the brute-force count of
survivors over one period equals $\prod_p(p-\nu_p)$ exactly (True for all
eight), and the frequency over $n\le10^7$ matches the product of local
frequencies to $7.11\times10^{-7}$ (bound $D/X=4.6\times10^{-3}$).
\end{observation}
\begin{observation}[Local volumes are Euler factors]
$\beta_p=(1-\nu_p/p)(1-1/p)^{-2}$ measures $2$ at $p=2$ and
$1-1/(p-1)^2$ at odd $p$ ($3/4,15/16,35/36,99/100$ for
$p=3,5,7,11$); partial products $\prod_{p\le P}\beta_p$ converge to
$2C_2=1.3203236$ with guard $1.2\times10^{-7}$ at $P=10^6$: the Tamagawa
reading reconstructs Hardy--Littlewood from local volumes alone.
\end{observation}

\section{Exact laws and the measured seam}
The architectural claims of Sections 5--6 are theorems: relative invariance
and the dilation character (Theorem 5.1; notebook v10, Addendum XVIII, T1);
CRT independence with explicit error $D/X$ (T2); local volumes
$\beta_p=1-1/(p-1)^2$ with Euler tail $O(1/(P\log P))$ (T3); and, at the
integer layer with fixed modulus, equidistribution of smooth numbers in
coprime progressions yields the cross-term extinction
$C(X)=1+O(1/\log X)$ (T4, after Tenenbaum, III.5). The prime layer ---
smoothness of $p-1$ along primes in congruence classes --- is stated as
Conjecture C1: measured at $10^6$ and $10^7$ with contracting deviations,
published guards and a classified interaction residual; it is an instance of
the Hardy--Littlewood family with smoothness constraints, and no proof of it
is claimed. The seam between T4 and C1 is itself measured (Addendum XVII).

\section{Scope, published refusals, reproducibility}
No open conjecture is proved or claimed. Two predictions were corrected by
measurement (naive Dickman model; a blind pooled model) and the corrections
are results; one self-detected coding error (a cumulative product) is
published as a refused run in the companion log. Companions:
\texttt{lissite\_corrigee.py}, \texttt{lissite\_u4\_10M.py},
\texttt{obstruction\_parite.py}, \texttt{groupe\_translations.py},
\texttt{tamagawa\_local.py}, \texttt{obstruction\_mod3.py},
\texttt{obstruction\_mod12.py}, \texttt{terme\_croise.py} --- pure Python, no
dependencies, archived on Zenodo with the notebook.

\begin{thebibliography}{9}
\bibitem{hl1923} G.H. Hardy and J.E. Littlewood, \emph{Some problems of
`Partitio Numerorum'}, Acta Math., 1923.
\bibitem{dickman1930} K. Dickman, \emph{On the frequency of numbers containing
prime factors of a certain relative magnitude}, Arkiv f\"or Mat., Astron. och
Fys., 1930.
\bibitem{granville2008} A. Granville, \emph{Smooth numbers: computational
number theory and beyond}, in Algorithmic Number Theory Lattices and Fields,
2008.
\bibitem{zhang2013} Y. Zhang, \emph{Bounded gaps between primes}, Annals of
Mathematics, 2013.
\bibitem{tenenbaum} G. Tenenbaum, \emph{Introduction to Analytic and
Probabilistic Number Theory}, Cambridge University Press, 1995 (Chap. III.5:
smooth numbers in arithmetic progressions).
\end{thebibliography}
\end{document}